\documentclass[11pt,a4paper]{article}

% ── Packages ──────────────────────────────────────────────────────────────────
\usepackage[T1]{fontenc}
\usepackage[utf8]{inputenc}
\usepackage{palatino}
\usepackage{microtype}
\usepackage{geometry}
\usepackage{amsmath,amssymb,amsthm}
\usepackage{mathtools}
\usepackage{physics}
\usepackage{booktabs}
\usepackage{array}
\usepackage{xcolor}
\usepackage{tcolorbox}
\usepackage{hyperref}
\usepackage{fancyhdr}
\usepackage{titlesec}
\usepackage{parskip}
\usepackage{enumitem}

% ── Page layout ───────────────────────────────────────────────────────────────
\geometry{
  a4paper,
  top=2.8cm, bottom=3.2cm,
  left=2.8cm, right=2.8cm,
  headheight=14pt
}

% ── Colours ───────────────────────────────────────────────────────────────────
\definecolor{deepblue}{HTML}{1a1a2e}
\definecolor{midblue}{HTML}{2d3a7a}
\definecolor{lightblue}{HTML}{eef2ff}
\definecolor{ruleblue}{HTML}{9999bb}
\definecolor{boxframe}{HTML}{2d3a7a}
\definecolor{notegrey}{HTML}{555566}
\definecolor{accentgold}{HTML}{7a5c00}

% ── Hyperref ──────────────────────────────────────────────────────────────────
\hypersetup{
  colorlinks=true,
  linkcolor=midblue,
  urlcolor=midblue,
  citecolor=midblue,
  pdftitle={Causal Self-Consistency of the Speed of Light in Fractal S3 Geometry},
  pdfauthor={Andrei Preece, Boris Batenin},
  pdfsubject={Conceptual consistency framework, fractal S3, causal boundary}
}

% ── Section formatting ────────────────────────────────────────────────────────
\titleformat{\section}
  {\large\bfseries\color{deepblue}}
  {\thesection.}{0.6em}{}
  [\vspace{1pt}{\color{ruleblue}\hrule height 0.6pt}\vspace{2pt}]

\titleformat{\subsection}
  {\normalsize\bfseries\color{midblue}}
  {\thesubsection}{0.6em}{}

\titlespacing{\section}{0pt}{18pt}{8pt}
\titlespacing{\subsection}{0pt}{11pt}{4pt}

% ── Header / footer ───────────────────────────────────────────────────────────
\pagestyle{fancy}
\fancyhf{}
\fancyhead[L]{\small\color{midblue}\textit{Causal Self-Consistency — Fractal $S^3$ Geometry}}
\fancyhead[R]{\small\color{midblue}Preece \& Batenin, 2026}
\fancyfoot[C]{\small\color{notegrey}\thepage}
\renewcommand{\headrulewidth}{0.5pt}
\renewcommand{\headrule}{\hbox to\headwidth{%
  \color{ruleblue}\leaders\hrule height \headrulewidth\hfill}}

% ── tcolorbox styles ──────────────────────────────────────────────────────────
\tcbuselibrary{skins, breakable}

\newtcolorbox{closurebox}{
  enhanced, breakable,
  colback=lightblue,
  colframe=boxframe,
  boxrule=0.9pt, arc=4pt,
  left=10pt, right=10pt, top=8pt, bottom=8pt,
}

\newtcolorbox{principlebox}[1][]{
  enhanced, breakable,
  colback=deepblue!5,
  colframe=deepblue,
  boxrule=1pt, arc=4pt,
  left=10pt, right=10pt, top=8pt, bottom=8pt,
  title={\small\bfseries Causal Self-Consistency Principle},
  fonttitle=\bfseries\color{white},
  coltitle=white,
  attach boxed title to top left={yshift=-2.5mm, xshift=7mm},
  boxed title style={colback=deepblue, arc=3pt, boxrule=0pt, left=4pt, right=4pt},
  #1
}

\newtcolorbox{notebox}{
  enhanced, breakable,
  colback=accentgold!6,
  colframe=accentgold!60,
  boxrule=0.6pt, arc=3pt,
  left=9pt, right=9pt, top=6pt, bottom=6pt,
  fontupper=\small\color{notegrey},
}

% ── Macros ────────────────────────────────────────────────────────────────────
\newcommand{\lP}{l_{\mathrm{P}}}
\newcommand{\Rzero}{R_0}
\newcommand{\Mzero}{M_0}
\newcommand{\ph}{\varphi}
\newcommand{\Vg}{V_{S^3}}

% ─────────────────────────────────────────────────────────────────────────────
\begin{document}
\thispagestyle{empty}
% ─────────────────────────────────────────────────────────────────────────────

% ── Title block ───────────────────────────────────────────────────────────────
\begin{center}
  {\LARGE\bfseries\color{deepblue}
    Causal Self-Consistency of the Speed of Light\\[5pt]
    in Fractal $S^3$ Geometry}\\[18pt]
  {\large Andrei Preece \quad \& \quad Boris Batenin}\\[7pt]
  {\small
    \href{mailto:andre.preece@web.de}{andre.preece@web.de}\quad
    ORCID: \href{https://orcid.org/0009-0004-5978-795X}{0009-0004-5978-795X}\\[2pt]
    \href{mailto:batenin@web.de}{batenin@web.de}\quad
    ORCID: \href{https://orcid.org/0009-0000-7027-8583}{0009-0000-7027-8583}\\[7pt]
    Berlin, 2026
  }
\end{center}

\vspace{5pt}
{\color{midblue}\hrule height 1.4pt}
\vspace{14pt}

% ── Abstract ──────────────────────────────────────────────────────────────────
\begin{center}
\begin{minipage}{0.93\linewidth}
\textbf{Abstract.}\;
This work does not derive the numerical value of the speed of light from independent
first principles. Instead, it formulates a \emph{causal self-consistency test} for
fractal $S^3$ volume--density evolution. The speed of light $c$ is taken as the
empirically observed invariant causal speed, and the question posed is whether the
proposed $S^3$ hierarchy admits a closed, self-consistent volume--density transition
under that invariant boundary condition. A golden-ratio scaling $R_n = R_0\ph^n$ is
used as a structural hierarchy, while any Planck-scale normalisation is treated only
as an illustrative calibration, not as a derivation. The resulting identity $c = c$
is interpreted as a \emph{closure condition}: it shows that the assumed geometric and
dynamical relations remain internally consistent under the empirical causal speed,
rather than providing an independent prediction of $c$.

\smallskip
\noindent\textbf{Keywords:} empirical invariant $c$;\;causal boundary;\;
self-consistency test;\;closure condition;\;fractal $S^3$ hierarchy;\;
structural ansatz;\;conceptual consistency framework.
\end{minipage}
\end{center}

\vspace{10pt}
{\color{ruleblue}\hrule}
\vspace{14pt}

% ─────────────────────────────────────────────────────────────────────────────
\section{Introduction}
% ─────────────────────────────────────────────────────────────────────────────

In standard physics the speed of light
\[
  c = 299\,792\,458\ \mathrm{m\,s^{-1}}
\]
is an empirically measured invariant, enshrined in the SI definition of the metre.
Its invariance is a postulate of special relativity, established through measurement
and not derived from a more primitive geometric principle.

The present work adopts a different, more modest aim. We \textbf{do not} claim to
derive $c$ from deeper constants or abstract first principles. Instead, we
\textbf{test whether the proposed fractal $S^3$ hierarchy is internally consistent
with the empirically observed invariant causal speed $c$}. Specifically:
\begin{enumerate}[leftmargin=2em, itemsep=3pt]
  \item $c$ is accepted as the empirically measured invariant causal speed. No
    attempt is made to predict its numerical value independently.
  \item The model examines whether the fractal $S^3$ volume--density transition
    admits a self-consistent \emph{fixed-point condition} when the causal boundary
    propagates at $c$.
  \item The closure identity $c = c$ that arises from substitution is interpreted
    as a fixed-point / closure condition, confirming internal geometric coherence
    rather than constituting a derivation.
\end{enumerate}

Because an observer inside a single causal universe cannot measure light from
``outside'' the universe, $c$ is operationally accessed only through internal
empirical procedures. The present model therefore treats $c$ as \emph{empirical
input} and tests whether the proposed $S^3$ structure closes consistently under
this input.

% ─────────────────────────────────────────────────────────────────────────────
\section{Fractal $S^3$ Hierarchy}
% ─────────────────────────────────────────────────────────────────────────────

\subsection{Golden-ratio scaling as structural ansatz}

The fractal $S^3$ manifold is parameterised by level radii
\[
  \boxed{R_n = \Rzero\,\ph^n,}
  \qquad n\in\mathbb{Z},
\]
where $\ph = \tfrac{1+\sqrt{5}}{2}\approx 1.618$ is the golden ratio and $n=0$
marks the recombination epoch, at which photons decouple from matter and propagate
freely.

\smallskip
\noindent\textbf{Status of the ansatz.}\;
The golden-ratio hierarchy is used as a minimal self-similar scaling structure.
Its mathematical motivation is the Perron--Frobenius dominance of the minimal
positive second-order recursion, but its physical realisation in the present
light-boundary model remains a \emph{structural ansatz}. We do not assert that
$\ph$ is the uniquely valid scaling ratio for an $S^3$ geometry; we adopt it as
a discrete self-similar hierarchy and test the resulting consistency.

\subsection{Optional Planck-scale calibration}

A Planck-scale seed may be used as a reference normalisation,
\[
  \Rzero = \lP\,\ph^N,
  \qquad \lP = \sqrt{\frac{\hbar G}{c^3}}\approx 1.616\times10^{-35}\ \mathrm{m},
\]
connecting the quantum scale to a cosmological boundary scale through the
golden-ratio hierarchy. For the recombination-horizon scale
$\Rzero\approx 2.536\times10^8\ \mathrm{m}$, this gives
\[
  N = \frac{\ln(\Rzero/\lP)}{\ln\ph}\approx 206.8 \;\Longrightarrow\; N = 207.
\]

\begin{notebox}
\textbf{Calibration note.}\quad
The value $N\approx 207$ is used only as one illustrative calibration connecting
a Planck-scale seed to a recombination-scale boundary. It is \emph{not} required
for the logical structure of the closure condition developed in
Section~\ref{sec:closure}, and it does not constitute a derivation of $c$ or of any
cosmological observable.
\end{notebox}

% ─────────────────────────────────────────────────────────────────────────────
\section{Volume, Density, and Mass Normalisation}
% ─────────────────────────────────────────────────────────────────────────────

\subsection{Global $S^3$ volume}

For a three-sphere $S^3$ of radius $R$, the exact global volume is
\[
  \Vg(R) = 2\pi^2 R^3.
\]
This is the geometrically correct definition for the compact, simply connected,
boundaryless $S^3$ manifold used in this framework.

\subsection{Local Euclidean approximation}

In the regime of small hyperspherical coordinate ($\chi\ll 1$), the $S^3$ volume
reduces locally to the Euclidean approximation
\[
  V_{\mathrm{local}}(R) \approx \tfrac{4}{3}\pi R^3.
\]

\begin{notebox}
\textbf{Scope of the approximation.}\quad
The local Euclidean form is employed below only for illustrative dynamical scaling
in the $\chi\ll 1$ regime. It does not replace the global $S^3$ volume definition,
and a fully covariant treatment would use $\Vg$ throughout.
\end{notebox}

\subsection{Mass as a closure normalisation parameter}

The matter density at the zeroth level is parameterised as
\[
  \rho_0 = \frac{\Mzero}{\tfrac{4}{3}\pi\Rzero^3}.
\]
The mass $\Mzero$ is a \emph{normalisation parameter chosen to test closure of the
dynamical relation under the empirical causal speed $c$}. It is not derived from
first principles, and its value is determined self-consistently at the closure step
(Section~\ref{sec:closure}).

% ─────────────────────────────────────────────────────────────────────────────
\section{Causal Boundary Dynamics}
% ─────────────────────────────────────────────────────────────────────────────

\subsection{Boundary velocity}

The rate at which the fractal boundary advances between consecutive levels is
\[
  c_n = \frac{R_{n+1}-R_n}{\Delta t_n}
       = \frac{\Rzero\ph^n(\ph-1)}{\Delta t_n}.
\]
At the zeroth level ($n=0$):
\[
  c = \frac{\Rzero(\ph-1)}{\Delta t_0}.
\]
We \emph{posit} that this boundary velocity equals the empirically measured $c$,
and subsequently test whether this assumption is maintained self-consistently.

\subsection{Effective expansion rate}

Using the Friedmann-type effective Hubble parameter,
\[
  H_0 = \sqrt{\frac{8\pi G\rho_0}{3}}
       = \sqrt{\frac{2G\Mzero}{\Rzero^3}},
\]
the causal boundary velocity becomes
\[
  c = H_0\,\Rzero = \sqrt{\frac{2G\Mzero}{\Rzero}}.
\]
Incorporating the fractal correction $1/\ph$ (arising from the golden-ratio step
between adjacent levels):
\begin{equation}\label{eq:cdyn}
  c = \sqrt{\frac{2G\Mzero}{\Rzero}}\cdot\frac{1}{\ph}.
\end{equation}

% ─────────────────────────────────────────────────────────────────────────────
\section{Closure Condition}
\label{sec:closure}
% ─────────────────────────────────────────────────────────────────────────────

Solving equation~\eqref{eq:cdyn} for the mass normalisation that makes the
relation consistent with the empirical invariant $c$ gives:
\[
  \Mzero = \frac{c^2\,\Rzero\,\ph^2}{2G}.
\]
Substituting back into~\eqref{eq:cdyn}:
\begin{align}
  c_{\mathrm{out}}
  &= \sqrt{\frac{2G}{\Rzero}\cdot\frac{c^2\Rzero\ph^2}{2G}}\cdot\frac{1}{\ph}
   = \sqrt{c^2\ph^2}\cdot\frac{1}{\ph}
   = c\ph\cdot\frac{1}{\ph}
   = c.\label{eq:closure}
\end{align}

\begin{closurebox}
\textbf{Closure condition.}\;
The substitution yields the identity $c = c$. This identity is expected by
construction once $\Mzero$ is chosen as the closure normalisation.

\smallskip
The point is \emph{not} to predict $c$ independently. The point is to show that
\textbf{the proposed dynamical chain admits a self-consistent solution when the
causal boundary speed is the empirical invariant $c$}: the assumed fractal $S^3$
scaling, density--volume relation, and causal-boundary dynamics form a closed,
coherent system under that boundary condition. The result is therefore a
\emph{fixed-point condition for causal self-consistency}, not an independent
prediction of the speed of light.
\end{closurebox}

\medskip
The structure of the constraint is the following. The golden-ratio hierarchy,
Planck-scale calibration, and Friedmann-type dynamics together admit a
self-consistent solution when the boundary propagation speed equals $c$. Modifying
the scaling ratio, the calibration level $N$, or the density--volume relation
changes the algebraic structure and may or may not admit a self-consistent closure.
Within the present model calibration, $c$ appears as the consistent propagation
speed — the \emph{fixed point of causal self-consistency}.

% ─────────────────────────────────────────────────────────────────────────────
\section{Analogy with Maxwell Electrodynamics}
% ─────────────────────────────────────────────────────────────────────────────

In Maxwell's electrodynamics, the speed of light appears as
\[
  c = \frac{1}{\sqrt{\varepsilon_0\,\mu_0}},
\]
relating the propagation speed of electromagnetic waves to the electric permittivity
$\varepsilon_0$ and magnetic permeability $\mu_0$ of the vacuum. This relation does
not derive $c$ ex nihilo. It relates $c$ to independently measured electromagnetic
constants and shows that the internal structure of Maxwell's equations singles out a
characteristic invariant speed.

Similarly, the present framework does not derive $c$ from nothing. It tests whether
a fractal $S^3$ volume--density hierarchy admits a causal closure when the boundary
speed is the empirical invariant $c$. In both cases, $c$ appears not as an arbitrary
free parameter, but as \emph{the speed required for internal consistency of the
theoretical structure}.

\medskip
\begin{notebox}
\textbf{Epistemological note.}\quad
This analogy is intended as an epistemological comparison only. We do not claim that
the present framework achieves the same empirical scope or predictive status as
Maxwell's theory. The analogy illustrates a shared logical structure: in both cases,
$c$ is singled out by an internal consistency requirement rather than introduced
as a free parameter.
\end{notebox}

% ─────────────────────────────────────────────────────────────────────────────
\section{Causal Self-Consistency Principle}
% ─────────────────────────────────────────────────────────────────────────────

\begin{principlebox}
A compact, simply connected, boundaryless $S^3$ geometry undergoing fractal
volume--density evolution governed by $R_n = R_0\ph^n$ admits a coherent
causal-boundary description when the boundary propagation speed is the empirical
invariant $c$. Within the present model calibration, the golden-ratio hierarchy,
Planck-scale normalisation, and Friedmann-type dynamics together select $c$ as the
consistent propagation speed — the fixed point of causal self-consistency.
\end{principlebox}

\medskip
This principle is understood as a \emph{necessary condition for internal coherence
within the stated model}. It is not asserted as a uniqueness theorem without further
derivation: asserting ``if and only if'' would require a separate proof that no
other boundary speed admits any reparametrisation of the model that also closes
consistently. The present paper establishes the fixed-point character of $c$, not
its uniqueness in a strong mathematical sense.

% ─────────────────────────────────────────────────────────────────────────────
\section{Numerical Illustration}
% ─────────────────────────────────────────────────────────────────────────────

We verify that the closure condition~\eqref{eq:closure} is numerically realised for
standard values of the physical constants, using the $N=207$ illustrative calibration.

\begin{center}
\renewcommand{\arraystretch}{1.35}
\begin{tabular}{llr}
  \toprule
  \textbf{Parameter} & \textbf{Symbol} & \textbf{Value} \\
  \midrule
  Planck length & $\lP$ & $1.616\times10^{-35}\ \mathrm{m}$ \\
  Golden ratio & $\ph$ & $1.618$ \\
  Base scale (calibration) & $\Rzero = \lP\ph^{207}$ & $2.536\times10^{8}\ \mathrm{m}$ \\
  Gravitational constant & $G$ & $6.674\times10^{-11}\ \mathrm{m^3\,kg^{-1}\,s^{-2}}$ \\
  Speed of light \emph{(empirical input)} & $c$ & $2.998\times10^{8}\ \mathrm{m\,s^{-1}}$ \\
  \bottomrule
\end{tabular}
\end{center}

\bigskip
\noindent\textbf{Closure normalisation:}
\[
  \Mzero = \frac{c^2\Rzero\ph^2}{2G}
  = \frac{(2.998\times10^8)^2\cdot 2.536\times10^8\cdot(1.618)^2}
         {2\times6.674\times10^{-11}}
  \approx 4.47\times10^{35}\ \mathrm{kg}.
\]

\noindent\textbf{Recovered causal boundary speed:}
\[
  c_{\mathrm{out}}
  = \sqrt{\frac{2G\Mzero}{\Rzero}}\cdot\frac{1}{\ph}
  = \sqrt{\frac{2\times6.674\times10^{-11}\times4.47\times10^{35}}
               {2.536\times10^8}}
    \cdot\frac{1}{1.618}
  \approx 2.998\times10^{8}\ \mathrm{m\,s^{-1}}.
\]

\begin{notebox}
This numerical agreement reflects the \emph{imposed closure condition} and does not
constitute an independent prediction of $c$. The recovered value equals the input
because $\Mzero$ was defined precisely to enforce this agreement. The numerical
exercise demonstrates that the closure is algebraically and numerically coherent,
not that $c$ has been derived from more primitive quantities.
\end{notebox}

% ─────────────────────────────────────────────────────────────────────────────
\section{Conclusion}
% ─────────────────────────────────────────────────────────────────────────────

The results of this paper may be summarised as follows.

\begin{enumerate}[leftmargin=2em, itemsep=5pt]

  \item \textbf{No derivation of $c$.}\;
  This work does not derive the numerical value of the speed of light from
  independent first principles. $c$ is treated throughout as the empirically
  observed invariant causal speed.

  \item \textbf{Closure structure.}\;
  The fractal $S^3$ hierarchy $R_n = R_0\ph^n$, the density--volume relation, and
  the causal-boundary dynamics together form a self-consistent closure structure
  when the boundary propagation speed is set to the empirical invariant $c$. The
  identity $c = c$ arising from substitution is a \emph{fixed-point condition for
  causal self-consistency}.

  \item \textbf{Epistemological analogy.}\;
  The logical structure parallels Maxwell's relation $c = 1/\!\sqrt{\varepsilon_0\mu_0}$:
  in both cases, $c$ is not an arbitrary free parameter but the speed required for
  internal consistency of the theoretical framework.

  \item \textbf{Optional calibration.}\;
  The Planck-scale normalisation $R_0 = \lP\ph^N$ with $N\approx207$ is an
  illustrative calibration connecting quantum and cosmological scales through the
  golden-ratio hierarchy. It is not a fundamental derivation and is not required for
  the logical structure of the closure condition.

  \item \textbf{Scope.}\;
  This paper is a conceptual consistency study in fractal $S^3$ geometry.
  Its contribution is to show that the proposed geometric and dynamical
  framework is internally coherent under the empirical causal speed $c$.

\end{enumerate}

\medskip
\begin{closurebox}
Thus, $c$ should not be understood here as a number derived from more primitive
constants, but as the \textbf{empirical invariant causal speed that makes the
proposed fractal $S^3$ volume--density evolution close consistently}. The result is
a fixed-point condition for causal self-consistency — a conceptual consistency
framework, not an independent prediction of the speed of light.
\end{closurebox}

\vspace{18pt}
{\color{ruleblue}\hrule}
\vspace{10pt}
\begin{center}
  \textit{\color{midblue}Namaste.}
\end{center}

\vspace{20pt}
{\color{midblue}\hrule height 1.2pt}
\vspace{12pt}

% ─────────────────────────────────────────────────────────────────────────────
\begin{thebibliography}{99}
\small

\bibitem{einstein1905}
A.~Einstein,
\textit{Zur Elektrodynamik bewegter K\"orper},
Annalen der Physik \textbf{17} (1905) 891--921.

\bibitem{maxwell1865}
J.~C.~Maxwell,
\textit{A Dynamical Theory of the Electromagnetic Field},
Philosophical Transactions of the Royal Society of London
\textbf{155} (1865) 459--512.

\bibitem{penrose2004}
R.~Penrose,
\textit{The Road to Reality: A Complete Guide to the Laws of the Universe},
Jonathan Cape, London, 2004.

\bibitem{misner1973}
C.~W.~Misner, K.~S.~Thorne, J.~A.~Wheeler,
\textit{Gravitation},
W.~H.~Freeman, San Francisco, 1973.

\bibitem{weinberg1972}
S.~Weinberg,
\textit{Gravitation and Cosmology: Principles and Applications
  of the General Theory of Relativity},
John Wiley \& Sons, New York, 1972.

\bibitem{kolb1990}
E.~W.~Kolb, M.~S.~Turner,
\textit{The Early Universe},
Addison-Wesley, Redwood City, 1990.

\bibitem{livio2002}
M.~Livio,
\textit{The Golden Ratio: The Story of Phi, the World's Most Astonishing Number},
Broadway Books, New York, 2002.

\bibitem{codata2018}
E.~R.~Cohen, T.~Cvita\v{s}, J.~G.~Frey \textit{et al.},
\textit{CODATA Internationally Recommended 2018 Values
  of the Fundamental Physical Constants},
NIST, Gaithersburg, MD, 2019.
\url{https://physics.nist.gov/cuu/Constants/}

\bibitem{planck1900}
M.~Planck,
\textit{\"Uber irreversible Strahlungsvorg\"ange},
Sitzungsberichte der K\"oniglich Preu\ss{}ischen Akademie der Wissenschaften
(1900) 440--480.

\bibitem{friedmann1922}
A.~Friedmann,
\textit{\"Uber die Kr\"ummung des Raumes},
Zeitschrift f\"ur Physik \textbf{10} (1922) 377--386.

\end{thebibliography}

\end{document}
